\documentclass{article}

\usepackage{amsmath, makecell}

\title{Ch. 30 Practice Quiz Solutions}
\author{Principles of Macroeconomics}
\date{\today}

\begin{document}

\maketitle

\begin{enumerate}

\item Stocks are a store of value, but they aren't used as a medium of exchange or a unit of account. You typically can't purchase goods with stocks, so they aren't a consistent medium of exchange. Prices aren't quoted in terms of stocks, so they aren't a unit of account.

\item The Fed's dual mandate is to keep inflation low and maintain maximum sustainable employment.

\item
	\begin{gather*}
	480 = c + d \\
	c = 2 \cdot res \\
	d = \frac{1}{0.1} \cdot res = 10 \cdot res\\
	480 = 2 \cdot res + 10 \cdot res \\
	40 = res \\
	c = 2 \cdot 40 = \$80 \\
	d = 10 \cdot 40 = \$400
	\end{gather*}
	
\item

	\begin{enumerate}
	
	\item
		\begin{tabular}{c | c}
		Assets & Liabilities \\
		\hline
		Reserves: $\$1,000 \cdot 0.1 = \$100$ & Deposits: \$1,000 \\
		Loans: \$1,100 & Debt: \$600 \\
		Securities: \$800 & \makecell{Capital: $\$100 + \$1,100 + \$800 -$ \\ $\$1,000 - \$600 = \$400$}
		\end{tabular}
		
	\item $\frac{2,000}{400} = 5$
	
	\item
		\begin{gather*}
		\$800(1.25) = \$1,000 \\
		\text{Assets} = \$100 + \$1,100 + \$1,000 = \$2,200 \\
		\% \Delta \text{Assets} = \frac{2,200 - 2,000}{2,000} = 10\% \\
		\text{Capital} = \$2,200 - \$1,000 - \$600 = \$600 \\
		\% \Delta \text{Capital} = \frac{600 - 400}{400} = 50\%
		\end{gather*}
		
	\end{enumerate}
	
\item

	\begin{enumerate}
	
	\item That increases the money supply because banks will hold fewer reserves, so the reserve ratio will go down, and the money multiplier will go up.
	
	\item That increases the money supply because banks will borrow less from the Fed and therefore hold fewer reserves. The will decrease the reserve ratio and increase the money multiplier.
	
	\item That decreases the money supply because the Fed will take the money it receives for the bonds out of circulation.
	
	\end{enumerate}

\end{enumerate}

\end{document}